
\chapter*{The Machine-Learning Map}
\addcontentsline{toc}{chapter}{The Machine-Learning Map}
\markboth{The Machine-Learning Map}{}

Describe a project by its feedback, task, data structure, model family, and lifecycle stage. Read the maps in that order.

\begin{keyidea}
Ask: \textbf{What output? What feedback? What data structure? How will future performance be measured?}
\end{keyidea}

\section*{1. The whole picture}

Machine learning connects AI with statistics, optimization, computing, and data systems. These overlapping views are more useful than a rigid hierarchy.

\begin{mlfigure}
\begin{tikzpicture}[
  box/.style={draw=MLBlue,fill=MLPaleBlue,rounded corners=1.8mm,
              text width=4.35cm,minimum height=1.5cm,align=center,
              font=\footnotesize\color{MLNavy}},
  core/.style={draw=MLGreen,fill=MLPaleTeal,rounded corners=2mm,
              minimum width=3.4cm,minimum height=1.0cm,align=center,
              font=\small\bfseries\color{MLNavy}},
  adj/.style={draw=MLMuted!70,fill=MLPaleGray,rounded corners=1.8mm,
              text width=4.35cm,minimum height=0.9cm,align=center,
              font=\scriptsize\color{MLNavy}},
  arr/.style={-{Stealth[length=2mm]},thick,draw=MLTeal}
]
\node[core] (ml) at (0,0) {MACHINE LEARNING\\\scriptsize\mdseries learn patterns from data or interaction};

\node[box] (feedback) at (-5.0,2.1) {\textbf{Learning signal}\\[2pt]labels or no labels;\\rewards from interaction};
\node[box] (task) at (0,2.1) {\textbf{Task}\\[2pt]predict, discover,\\generate, decide};
\node[box] (data) at (5.0,2.1) {\textbf{Data structure}\\[2pt]tables, time, text,\\images, graphs};

\node[box] (family) at (-5.0,-2.1) {\textbf{Model family}\\[2pt]linear models, trees,\\kernels, neural networks};
\node[box] (eval) at (0,-2.1) {\textbf{Evaluation}\\[2pt]splits, metrics,\\uncertainty, error analysis};
\node[box] (life) at (5.0,-2.1) {\textbf{Lifecycle}\\[2pt]train and deploy,\\monitor and govern};

\draw[arr] (feedback) -- (ml);
\draw[arr] (task) -- (ml);
\draw[arr] (data) -- (ml);
\draw[arr] (ml) -- (family);
\draw[arr] (ml) -- (eval);
\draw[arr] (ml) -- (life);

\node[adj] at (-5.0,3.65) {Statistics and probability};
\node[adj] at (0,3.65) {Optimization and linear algebra};
\node[adj] at (5.0,3.65) {Computer and data systems};
\end{tikzpicture}
\end{mlfigure}

\begin{mltable}
\begin{tabularx}{\linewidth}{@{}>{\bfseries\leavevmode\color{MLNavy}}p{0.18\linewidth}X@{}}
\toprule
\rowcolor{MLPaleBlue!92}
View & Main branches \\
\midrule
Feedback & Full, partial, or no external labels; interaction rewards. \\
Task & Predict, rank, forecast, discover structure, represent, generate, detect anomalies, or choose actions. \\
Data & Tables, time series, text, images/video, audio, graphs, spatial, or multimodal data. \\
Model family & Linear, probabilistic, neighbor, tree/ensemble, kernel, neural, generative, or graph models. \\
Evaluation & Realistic splits, baselines, metrics, calibration, uncertainty, and subgroup errors. \\
Lifecycle & Frame, collect, prepare, train, evaluate, deploy, monitor, maintain, govern. \\
\bottomrule
\end{tabularx}
\end{mltable}

\section*{2. Start from the feedback you actually have}

\begin{mlfigure}
\begin{tikzpicture}[
  q/.style={draw=MLNavy,fill=MLNavy,text=white,rounded corners=2mm,
            text width=14.05cm,minimum height=0.9cm,align=center,font=\bfseries\small},
  branch/.style={draw=MLBlue,fill=MLPaleBlue,rounded corners=2mm,
                 text width=3.25cm,minimum height=1.0cm,align=center,font=\small\color{MLNavy}},
  leaf/.style={draw=MLTeal,fill=MLPaleTeal,rounded corners=2mm,
               text width=3.0cm,minimum height=1.58cm,align=center,font=\scriptsize\color{MLNavy}},
  arr/.style={-{Stealth[length=2mm]},thick,draw=MLMuted}
]
\node[q] (start) at (0,3.8) {What feedback is available?};

\node[branch] (full) at (-5.4,1.9) {A label for each example};
\node[branch] (few) at (-1.8,1.9) {Only some labels};
\node[branch] (none) at (1.8,1.9) {No external label};
\node[branch] (reward) at (5.4,1.9) {Reward after actions};

\node[leaf] (sup) at (-5.4,-0.1) {\textbf{Supervised learning}\\regression, classification, ranking, forecasting, survival};
\node[leaf] (limited) at (-1.8,-0.1) {\textbf{Limited-label learning}\\semi-supervised, active learning};
\node[leaf] (unsup) at (1.8,-0.1) {\textbf{Unsupervised / self-supervised}\\clustering, representation, density, generation};
\node[leaf] (rl) at (5.4,-0.1) {\textbf{Reinforcement learning}\\bandits, policies, value learning, control};

\draw[arr] (start.south) -- (0,2.9);
\draw[arr] (-5.4,2.9) -- (-5.4,2.4);
\draw[arr] (-1.8,2.9) -- (-1.8,2.4);
\draw[arr] (1.8,2.9) -- (1.8,2.4);
\draw[arr] (5.4,2.9) -- (5.4,2.4);
\draw[draw=MLMuted,thick] (-5.4,2.9) -- (5.4,2.9);

\draw[arr] (full) -- (sup);
\draw[arr] (few) -- (limited);
\draw[arr] (none) -- (unsup);
\draw[arr] (reward) -- (rl);
\end{tikzpicture}
\end{mlfigure}

\begin{intuition}
\textbf{Self-supervised learning} derives targets from data, such as hidden words. \textbf{Active learning} instead chooses examples for human labeling.
\end{intuition}

\section*{3. Start from the job you need done}

\begin{mltable}
\begin{tabularx}{\linewidth}{@{}>{\bfseries\leavevmode\color{MLNavy}}p{0.30\linewidth}p{0.30\linewidth}X@{}}
\toprule
\rowcolor{MLPaleBlue!92}
Need & First framing & Where it lives \\
\midrule
Predict a number & Regression & Core: \chapref{ch:regression} \\
Predict a category & Classification & Core: \chapref{ch:classification} \\
Predict an ordered category & Ordinal classification / regression & Introduced in \chapref{ch:classification} \\
Rank a finite set of items & Ranking / recommendation & Beyond the core; evaluate at the operational cutoff \\
Predict later values & Forecasting / temporal prediction & Uses supervised learning with chronological validation \\
Estimate time until an event & Survival / event-time modeling & Beyond the core; censoring changes the evaluation problem \\
Find groups & Clustering & Core: \chapref{ch:unsupervised}, then \chapref{ch:clustering} \\
Compress or build a representation & PCA, embeddings, representation learning & Core PCA: \chapref{ch:dimred}; modern embeddings are beyond \\
Find unusual cases & Anomaly detection & Core: \chapref{ch:anomaly} \\
Model a full distribution or generate samples & Density estimation / generative modeling & GMM is introduced in \chapref{ch:clustering}; modern generative models are beyond \\
Choose actions over time & Bandits / reinforcement learning & Extension: \chapref{ch:rl} \\
Estimate what an intervention will change & Causal inference / experimentation & Beyond the core; requires causal assumptions or experiments \\
\bottomrule
\end{tabularx}
\end{mltable}

\begin{keyidea}
Prediction estimates outcomes; causal inference estimates intervention effects. Predictive success does not establish causality.
\end{keyidea}

\section*{4. Start from the data structure}

The input representation often narrows the sensible model families before tuning begins.

\begin{mltable}
\renewcommand{\arraystretch}{1.12}
\setlength{\tabcolsep}{4.5pt}
\footnotesize
\begin{tabularx}{\linewidth}{@{}>{\raggedright\arraybackslash\bfseries\leavevmode\color{MLNavy}}p{0.21\linewidth}>{\raggedright\arraybackslash}p{0.33\linewidth}>{\raggedright\arraybackslash}X@{}}
\toprule
\rowcolor{MLPaleBlue!92}
\bfseries Data & \bfseries Strong classical starting point & \bfseries Modern directions and the main caution \\
\midrule
Tabular rows & Linear/logistic models, trees, forests, boosting & Deep tabular models; compare with strong tree baselines. \\
Time series / events & Lagged features with linear/tree models & Sequence models, transformers; respect chronology, horizon, and label timing. \\
Text & Counts/TF--IDF with naive Bayes, logistic regression, or SVM & Embeddings, language models; assess meaning, not just fluency. \\
Images / video & Handcrafted features, PCA, linear/SVM models & CNNs, vision transformers, diffusion; consider transfer learning. \\
Audio / speech & Spectral features with classical classifiers & Audio encoders, transformers; check sampling and alignment. \\
Graphs / networks & Node, edge, and graph statistics & Graph neural networks; avoid leakage through connections. \\
Spatial data & Domain features with linear/tree models & Spatial statistics, geospatial networks; account for nearby dependencies. \\
Multimodal & Combine separate feature pipelines & Multimodal models; align modalities carefully. \\
\bottomrule
\end{tabularx}
\end{mltable}

\section*{5. Model families are reusable building blocks}

Model families are reusable across tasks: trees, kernels, and networks support multiple prediction problems.

\enlargethispage{4\baselineskip}
\begin{mltable}
\renewcommand{\arraystretch}{1.12}
\setlength{\tabcolsep}{4.5pt}
\footnotesize
\begin{tabularx}{\linewidth}{@{}>{\raggedright\arraybackslash\bfseries\leavevmode\color{MLNavy}}p{0.23\linewidth}>{\raggedright\arraybackslash}p{0.30\linewidth}>{\raggedright\arraybackslash}X@{}}
\toprule
\rowcolor{MLPaleBlue!92}
\bfseries Family & \bfseries Core idea & \bfseries Examples in or beyond these notes \\
\midrule
Baselines & Simple reference rule & Mean, median, majority, random, or domain rule \\
Linear / GLMs & Weighted feature combinations & Linear, ridge, lasso, logistic, softmax \\
Probabilistic & Model probability distributions & Naive Bayes, LDA, GMM; Bayesian models beyond the core \\
Neighbors & Use nearby cases & $k$-NN, local outlier methods, density neighborhoods \\
Trees / ensembles & Partition space; combine learners & CART, random forests, boosting, stacking \\
Kernels / margins & Separate in a transformed space & SVM, SVR, one-class SVM \\
Neural networks & Learn features and predictions together & MLP, CNN, transformers, graph neural networks \\
Generative / foundation & Learn distributions or reusable representations & Autoencoders, diffusion, language and multimodal models \\
\bottomrule
\end{tabularx}
\end{mltable}

\begin{keyidea}
Compare a baseline, simple model, and stronger model in turn. Stop when further gains no longer justify complexity.
\end{keyidea}

\section*{6. Good first models for the core problems}

\begin{mltable}
\begin{tabularx}{\linewidth}{@{}>{\bfseries\leavevmode\color{MLNavy}}p{0.28\linewidth}X X@{}}
\toprule
\rowcolor{MLPaleBlue!92}
Problem & Start with & Then compare with \\
\midrule
Regression & Mean baseline, linear regression & Ridge/lasso, trees, random forest, gradient boosting, support vector regression \\
Binary classification & Majority baseline, logistic regression & $k$-NN, naive Bayes, trees, random forest, boosting, SVM \\
Multiclass classification & Majority baseline, softmax regression & Trees, random forest, boosting, $k$-NN, SVM \\
Clustering & $K$-means after scaling & Hierarchical clustering, DBSCAN, Gaussian mixtures \\
Dimensionality reduction & PCA & t-SNE or UMAP for exploratory visualization \\
Anomaly detection & Simple statistical rules & Local methods, isolation forest, one-class models \\
\bottomrule
\end{tabularx}
\end{mltable}

\section*{7. What these notes cover --- and what lies beyond}

\begin{mltable}
\begin{tabularx}{\linewidth}{@{}>{\bfseries\leavevmode}p{0.22\linewidth}X@{}}
\toprule
\rowcolor{MLPaleBlue!92}
\color{MLNavy}Scope & \bfseries\color{MLNavy}Topics \\
\midrule
\color{MLGreen}Core in these notes & \normalfont Problem framing and leakage-safe preparation; baselines and evaluation; classical supervised and unsupervised methods. \\
\color{MLBlue}Introduced / bridged & \normalfont Limited-label and reinforcement learning; neural networks, CNNs, and LSTMs on small data; transfer learning; deployment, drift, and responsibility. \\
\color{MLOrange!85!black}Beyond the core & \normalfont Deep, foundation, generative, graph, Bayesian, causal, online, federated, privacy-preserving, and production ML. \\
\bottomrule
\end{tabularx}
\end{mltable}
